\documentclass[12pt]{article}
\usepackage[margin=1in]{geometry}
\usepackage{setspace}
\RequirePackage[colorlinks,citecolor=blue,linkcolor=blue,urlcolor=blue]{hyperref}
\doublespacing

\begin{document}

\begin{center}
    \singlespacing
    {\Large\bf Identifying Rational Types in Unknown Environments under an Indirect Dynamic Mechanism}
    
    \vspace{1.5cm}
    
    {\large Humberto Bernal\footnote{Economist, professor of the economics program at the Universidad Colegio Mayor de Cundinamarca and Ph.D. in economics from the Universidad de los Andes, Colombia. E-mail: \href{mailto:hbernal@uniandes.edu.co}{hbernal@uniandes.edu.co}, \href{mailto:hbernalc@universidadmayor.edu.co}{hbernalc@universidadmayor.edu.co}. ORCID: \href{https://orcid.org/0000-0002-4864-1726}{0000-0002-4864-1726}. All research is the sole responsibility of the author, received no funding, and was developed as an independent initiative.}}
    
    \vspace{0.5cm}
    
    {\small Economics Program, Universidad Colegio Mayor de Cundinamarca; Ph.D. in Economics, Universidad de los Andes, Colombia.}
    
    \vspace{1.5cm}
    
    \textbf{Abstract}
\end{center}

\begin{abstract}
This paper develops a dynamic indirect mechanism to identify rational types in unknown environments characterized by observational pooling. The central theoretical contribution introduces the Mano de Dios--Guadalupe (MDG) process, an endogenous implementation noise kernel that maps latent intentions into executed actions with full support in the Bayesian game. Integrating the MDG process with a competing risks framework, the Rational Type Identification Theorem is established, whereby posterior Bayesian beliefs converge almost surely to the true type. Furthermore, the Implementability Theorem is presented, which guarantees that the mechanism is executable in a perfect Bayesian equilibrium. The empirical application to kidnapping for ransom in Colombia validates the theoretical framework through a microeconometric calibration and structural simulations that demonstrate the robustness of both Bayesian learning and incentive compatibility constraint compliance.
\end{abstract}

\vspace{1cm}

\noindent \textbf{Keywords:} indirect mechanism design, dynamic screening, endogenous implementation noise, Bayesian learning, type identification \\
\noindent \textbf{JEL Codes:} C73, D82, K42

\end{document}
